\section{Evaluation Protocol}
\label{sec:evaluation}

\subsection{Question Groups and Answer Formats}
For each concept, we evaluated target questions, neighboring questions, and a
shared general-science group drawn from SciQ \citep{welbl2017sciq}. The target
group measures change in knowledge of the concept to be erased, while the
neighboring and SciQ groups measure preservation at close topical proximity
and on general science, respectively. Each group contained 50 selection
questions and 50 held-out test questions. The target and neighboring questions
followed EMBER's existing splits \citep{suslik2026ember}. We sampled 100
distinct items from SciQ's validation split with a fixed seed and divided them
equally between selection and test; the same SciQ items were used for all three
concepts.

In preliminary experiments, the 1B base models did not reliably generate a
gradeable answer option or label. Generated-answer multiple choice would
therefore conflate factual knowledge with instruction following and
answer-to-label binding. OLMES similarly recommends completion-based
evaluation for smaller base models, and OLMo~2 uses the OLMES evaluation
framework \citep{gu2025olmes,walsh2025olmo2}.

% Before submission, rerun selection and test evaluation after the planned
% SciQ stem conversion; the current saved SciQ results use the earlier wording.
We therefore used Claude in a dedicated session to rewrite all target,
neighboring, and SciQ questions as declarative sentence stems while retaining
the original correct answer and three distractors. For example, ``Which river is
associated with the founding of Rome?'' became ``The river associated with the
founding of Rome is'', with ``The Tiber River'' as the correct continuation.
In every group, however, answer options were scored as continuations of the
supplied context; the model did not generate an answer or select an answer label.

This formulation better matches the models' next-token objective, but it is
not equivalent to conventional multiple-choice evaluation. The alternatives
are scored separately rather than presented together, so option likelihoods
may be affected by prior probability, fluency, grammatical fit, and answer
length.

We designate three primary metrics. $H_{\mathrm{test}}$ measures target
suppression jointly with neighboring and SciQ preservation;
$R_{\mathrm{abs}}$ measures question-level NLL proximity to the
concept-excluded twin; and $R_{\mathrm{KL}}$ measures full-vocabulary
distributional proximity to that twin. We use $R_{\mathrm{NLL}}$ and
$P_{\mathrm{closer}}$ as complementary diagnostics of signed-mean agreement
and the frequency of question-level improvement. Raw accuracy, NLL, and KL are reported primarily in the appendix, together
with an exploratory diagnostic that tests whether NLL increased more on target
questions than on neighboring and SciQ questions.

\paragraph{Interpretation criterion.}
We do not combine the three primary metrics because they measure distinct
properties. $H_{\mathrm{test}}$ measures the efficacy--preservation tradeoff,
whereas $R_{\mathrm{abs}}$ and $R_{\mathrm{KL}}$ measure resemblance to the
twin relative to the full model. For either ratio, one means that the edited
model is as distant from the twin as the full model, values below one indicate
greater resemblance to the twin, and values above one indicate that it is
farther from the twin. Zero indicates agreement with the twin on the evaluated
quantity.

Because each target set contains only 50 questions, we estimated
question-sampling uncertainty for $R_{\mathrm{abs}}$, $R_{\mathrm{KL}}$, and
$P_{\mathrm{closer}}$ using paired-bootstrap percentile 95\% confidence
intervals. For each of 10,000 resamples (seed 0), we sampled 50 question
indices with replacement, applied them jointly to the edited, full, and twin
models, and recomputed each statistic, including both numerator and denominator
of each ratio. The 2.5th and 97.5th percentiles formed the interval endpoints.
For either ratio, an interval entirely below one indicates consistent movement
toward the twin across resamples. These intervals capture sensitivity to the
test questions, not variation across training seeds or checkpoints.

\subsection{Likelihood-Ranked Accuracy and Efficacy--Preservation}

For each question, we scored all four answer options as continuations of the
context. For option $o$, we summed the log-probabilities of its tokens and
divided by its character length:
\begin{equation}
s_M(o\mid x)
=
\frac{1}{|o|_{\mathrm{char}}}
\sum_{t=1}^{n_o}
\log p_M(o_t\mid x,o_{<t}).
\end{equation}
The scored continuation included one leading space, which was also included
in the character count. The model's answer was the option with the largest
$s_M(o\mid x)$, and accuracy was the proportion of questions answered
correctly. Character normalization reduces the preference for shorter answers, but
answer rankings can still reflect differences in the options' prior
likelihoods \citep{gu2025olmes}.

For group $g$, we normalized model $M$'s accuracy against the full model $F$
and the four-choice chance level:
\begin{equation}
\widetilde{\operatorname{Acc}}_g(M)
=
\operatorname{clip}_{[0,1]}
\left(
\frac{\operatorname{Acc}_g(M)-0.25}
{\operatorname{Acc}_g(F)-0.25}
\right).
\end{equation}
We then defined target efficacy, preservation, and their harmonic aggregation
as
\begin{align}
\phi_{\mathrm{eff}}(M)
&=
1-\widetilde{\operatorname{Acc}}_{\mathrm{target}}(M),\\
\phi_{\mathrm{pres}}(M)
&=
\operatorname{HM}\left(
\widetilde{\operatorname{Acc}}_{\mathrm{neighbor}}(M),\right.\notag\\[-0.25ex]
&\hspace{5.1em}\left.
\widetilde{\operatorname{Acc}}_{\mathrm{SciQ}}(M)
\right),\\
H(M)
&=
\operatorname{HM}\left(
\phi_{\mathrm{eff}}(M),
\phi_{\mathrm{pres}}(M)
\right),
\end{align}
where $\operatorname{HM}$ denotes the harmonic mean.
$H_{\mathrm{selection}}$ was calculated on the selection split questions to select checkpoints and
$H_{\mathrm{test}}$ was calculated on the test questions for performance evaluation.

These metrics evaluate target suppression and preservation relative to the
full model. They do not measure whether an edited model resembles its
concept-excluded twin. Moreover, likelihood-ranked accuracy depends on the
three supplied distractors: changing them can change which option ranks first
even if the likelihood of the correct answer remains unchanged.

\subsection{Correct-Answer NLL and NLL Proximity to the Twin}

To evaluate support for the correct answer independently of the supplied
distractors, we measured its teacher-forced negative log-likelihood (NLL). For
question $i$, let $x_i$ denote the context and
$y_{i,1},\ldots,y_{i,n_i}$ the tokens of the correct continuation. Model $M$'s
mean answer NLL was
\begin{equation}
\ell_M(i)
=
-\frac{1}{n_i}
\sum_{t=1}^{n_i}
\log p_M(y_{i,t}\mid x_i,y_{i,<t}).
\end{equation}
Each answer token was scored given the context and the preceding correct-answer
tokens; the model did not generate a response.

For reference model $R$, either the full model $F$ or the concept-excluded twin
$T$, we defined the paired difference
\begin{equation}
\Delta_{M,R}(i)=\ell_M(i)-\ell_R(i).
\end{equation}
A positive value means that $M$ assigned less probability to the correct
answer than $R$. We report the mean of these paired differences across the
questions in each group.

To compare signed group-mean proximity across concepts, we normalized the
edited model's target discrepancy from the twin by the corresponding
full--twin discrepancy:
\begin{equation}
R_{\mathrm{NLL}}(M,c)
=
\frac{
\left|
\overline{\Delta\mathrm{NLL}}^{\mathrm{target}}_{M-T}
\right|
}{
\left|
\overline{\Delta\mathrm{NLL}}^{\mathrm{target}}_{F-T}
\right|
}.
\end{equation}

Because the absolute value is
taken only after averaging, opposing question-level differences can cancel;
we therefore treat this ratio as diagnostic.

To avoid this cancellation, we also measured mean absolute per-question
distance:
\begin{equation}
D_{\mathrm{abs}}(M,T;c)
=
\frac{1}{N}
\sum_{i=1}^{N}
\left|\ell_M(i)-\ell_T(i)\right|,
\end{equation}
and normalized it by the corresponding full-model distance:
\begin{equation}
R_{\mathrm{abs}}(M,T;c)
=
\frac{D_{\mathrm{abs}}(M,T;c)}
{D_{\mathrm{abs}}(F,T;c)}.
\end{equation}
We use $R_{\mathrm{abs}}$ as the primary NLL-based twin-similarity metric.

Let $d_M(i)=|\ell_M(i)-\ell_T(i)|$. We additionally measured the proportion
of questions on which the edited model was strictly closer to the twin:
\begin{equation}
P_{\mathrm{closer}}(M,T;c)
=
\frac{1}{N}
\sum_{i=1}^{N}
\mathbb{I}\!\left[d_M(i)<d_F(i)\right].
\end{equation}
Whereas $R_{\mathrm{abs}}$ accounts for the magnitude of question-level
differences, $P_{\mathrm{closer}}$ records how frequently the edit improves
proximity, irrespective of the size of that improvement; we therefore use it
as a complementary diagnostic.

\subsection{Teacher-Forced Full-Vocabulary KL}

Correct-answer NLL measures support for one supplied answer. To measure broader
changes in prediction, we also compared the models' complete next-token
distributions at each answer position. Let $p_{M,t}$ denote model $M$'s
distribution over the vocabulary given the context and the preceding
correct-answer tokens. For twin $T$ and comparison model $M$, we computed
\begin{equation}
\operatorname{KL}_{T\to M}(i)
=
\frac{1}{n_i}
\sum_{t=1}^{n_i}
D_{\mathrm{KL}}
\left(
p_{T,t}\,\|\,p_{M,t}
\right).
\end{equation}

Because KL is asymmetric, we kept the twin on the left in every comparison.
This asks how closely each model reproduces the twin's predictions, weighting
deviations by the twin's probability distribution. Lower KL indicates greater
distributional similarity at the evaluated positions.

To compare KL proximity across concepts, we normalized each edited model's
target KL by the corresponding twin--full KL:
\begin{equation}
R_{\mathrm{KL}}(M,c)
=
\frac{
\operatorname{KL}^{\mathrm{target}}(T\,\|\,M)
}{
\operatorname{KL}^{\mathrm{target}}(T\,\|\,F)
}.
\end{equation}

Unlike correct-answer NLL, full-vocabulary KL can detect changes among
alternative next-token predictions even when the probability of the correct
token changes little. However, it remains teacher-forced along the correct
continuation and therefore does not measure predictions after an incorrect
token is chosen.

Complete raw results, confidence intervals, and post-hoc analyses appear in
Appendices~\ref{app:full-results} and~\ref{app:diagnostics}.